\section{Flight Routes Check}
\label{sec:cses-flight-routes-check}

\problemstatement{Given a directed graph of $n$ cities and $m$ flights,
determine whether you can travel between \emph{every} pair of cities. If so
output \texttt{YES}; otherwise output two cities $a$, $b$ such that there is
no route from $a$ to $b$.}

\begin{patternbox}
Strong connectivity can be checked without full SCC decomposition: a graph
is strongly connected iff, from any single node $s$, \emph{everyone is
reachable} in the original graph \emph{and} \emph{everyone can reach $s$}
(reachable in the reversed graph). Two DFS runs suffice.
\end{patternbox}

\subsection*{Forward and reverse reachability from one node}

Pick $s=1$. DFS in the original graph: if some node $u$ is unreachable, then
there is no route $1\to u$, so print ``$1\ u$''. Otherwise DFS in the graph
with all edges reversed: if some node $u$ is unreachable there, it cannot
reach $1$ in the original, so print ``$u\ 1$''. If both DFS runs cover every
node, then $1$ can reach all and all can reach $1$, which makes every pair
connected via $1$ -- output \texttt{YES}.

\begin{ideabox}[Strong Connectivity = Reach All and Be Reached from One Node]
If a single node both reaches everyone and is reached by everyone, any two
nodes $a,b$ connect through it ($a\to s\to b$). So two reachability
searches -- one forward, one reversed -- decide strong connectivity and
produce a counterexample pair when it fails.
\end{ideabox}

\subsection*{Algorithm}

\begin{enumerate}
  \item DFS/BFS from $1$ in the original graph; if node $u$ unreached,
        print ``$1\ u$''.
  \item DFS/BFS from $1$ in the reversed graph; if node $u$ unreached,
        print ``$u\ 1$''.
  \item Else print \texttt{YES}.
\end{enumerate}

\subsection*{Dry run}

\dryrun{$1\!\to\!2$, $2\!\to\!3$ (no way back); $n=3$.}

\begin{itemize}
  \item Reversed DFS from $1$ cannot reach $2$ (edges reversed:
        $2\!\to\!1$, $3\!\to\!2$), so $2$ cannot reach $1$.
  \item Output ``$2\ 1$''.
\end{itemize}

\subsection*{C++ implementation}

\begin{lstlisting}
#include <bits/stdc++.h>
using namespace std;

void bfs(int s, vector<vector<int>>& g, vector<bool>& vis) {
    queue<int> q; q.push(s); vis[s] = true;
    while (!q.empty()) {
        int u = q.front(); q.pop();
        for (int v : g[u]) if (!vis[v]) { vis[v] = true; q.push(v); }
    }
}

int main() {
    int n, m;
    cin >> n >> m;
    vector<vector<int>> adj(n + 1), radj(n + 1);
    for (int i = 0; i < m; i++) {
        int a, b; cin >> a >> b;
        adj[a].push_back(b);
        radj[b].push_back(a);
    }

    vector<bool> vis(n + 1, false);
    bfs(1, adj, vis);
    for (int u = 1; u <= n; u++)
        if (!vis[u]) { cout << 1 << " " << u << "\n"; return 0; }

    fill(vis.begin(), vis.end(), false);
    bfs(1, radj, vis);
    for (int u = 1; u <= n; u++)
        if (!vis[u]) { cout << u << " " << 1 << "\n"; return 0; }

    cout << "YES\n";
    return 0;
}
\end{lstlisting}

\begin{complexitybox}
\textbf{Time Complexity.} $O(n+m)$ -- two traversals. \textbf{Space
Complexity.} $O(n+m)$ for both adjacency lists.
\end{complexitybox}

\begin{takeawaybox}
Strong connectivity needs only two searches from one node -- forward and on
the reversed graph -- not a full SCC pass. The unreached node in either
search is a ready-made counterexample pair, which is exactly what the
problem asks for.
\end{takeawaybox}
